% Part: history 
% Chapter: biographies 
% Section: julia-robinson
\documentclass[../../../include/open-logic-section]{subfiles}

\begin{document}

\olfileid[hi]{his}{bio}{rob}

\olsection{जूलिया रॉबिन्सन}

\olphoto{robinson-julia}{जूलिया रॉबिन्सन}

जूलिया बोमन रॉबिन्सन अमेरिकी गणितज्ञ थीं। वे मुख्यतः निर्णय समस्याओं पर अपने
कार्य और सर्वाधिक प्रसिद्ध रूप से हिल्बर्ट की दसवीं समस्या के समाधान में
अपने योगदान के लिए जानी जाती हैं। रॉबिन्सन का जन्म 8 दिसंबर 1919 को
मिसूरी के सेंट~लुई में हुआ। रॉबिन्सन याद करती हैं कि बाल्यावस्था में ही
संख्याएँ उन्हें आकर्षित करती थीं \citep[4]{Reid1986}। नौ वर्ष की आयु में
उन्हें लाल बुखार हुआ और फिर आमवाती ज्वर के कई बार लौटने वाले दौरे पड़े। इससे
उन्हें बहुत समय बिस्तर पर बिताना पड़ा और वे पढ़ाई में पीछे रह गईं। यद्यपि
निजी शिक्षकों की सहायता से वे बराबरी पर आ सकीं, उनके रोग के शारीरिक प्रभावों
का उनके जीवन पर स्थायी असर रहा।

बचपन के संघर्षों के बावजूद रॉबिन्सन ने गणित और विज्ञानों के कई पुरस्कारों
के साथ माध्यमिक विद्यालय पूरा किया। उन्होंने सैन डिएगो स्टेट कॉलेज से
विश्वविद्यालयी जीवन आरंभ किया और अंतिम वर्ष में यूनिवर्सिटी ऑफ कैलिफोर्निया,
बर्कली चली गईं। वहाँ गणितज्ञ राफ़ाएल रॉबिन्सन ने उन्हें प्रभावित किया। दोनों
अच्छे मित्र बने और 1941 में विवाह किया। संकाय सदस्य की पत्नी होने के कारण
रॉबिन्सन को बर्कली के गणित विभाग में पढ़ाने से रोक दिया गया। यद्यपि वे गणित
की कक्षाएँ सुनती रहीं, उन्हें विश्वविद्यालय छोड़कर परिवार आरंभ करने की आशा
थी। किंतु विवाह के कुछ ही समय बाद रॉबिन्सन को निमोनिया हो गया। उन्हें बताया
गया कि बचपन के आमवाती ज्वर के कारण उनके हृदय पर बहुत-सा घाव-ऊतक जमा हो गया
था। इसकी गंभीरता के कारण चिकित्सक ने अनुमान लगाया कि वे चालीस वर्ष से अधिक
जीवित नहीं रहेंगी और उन्हें संतान न करने की सलाह दी गई
\citep[13]{Reid1986}।

रॉबिन्सन लंबे समय तक अवसादग्रस्त रहीं, किंतु अंततः उन्होंने गणित का अध्ययन
जारी रखने का निर्णय लिया। वे बर्कली लौटीं और 1948 में अल्फ्रेड टार्स्की के
निर्देशन में पीएचडी पूरी की। टार्स्की दिखा चुके थे कि वास्तविक संख्याओं का
प्रथम-क्रम सिद्धांत निर्णेय है, और ग्योडेल के कार्य से निष्पन्न होता था कि
प्राकृतिक संख्याओं का प्रथम-क्रम सिद्धांत अनिर्णेय है। परिमेय संख्याओं का
प्रथम-क्रम सिद्धांत निर्णेय है या नहीं, यह एक बड़ी खुली समस्या थी। अपने
शोधप्रबंध \citeyearpar{Robinson1949} में रॉबिन्सन ने सिद्ध किया कि वह निर्णेय
नहीं है।

निर्णय समस्याओं में रुचि के कारण रॉबिन्सन ने आगे हिल्बर्ट की दसवीं समस्या
का समाधान खोजने का प्रयास किया। यह समस्या 1900 में डेविड हिल्बर्ट द्वारा
रखी गई 23 प्रसिद्ध गणितीय समस्याओं की सूची में से एक थी। दसवीं समस्या पूछती
है कि क्या कोई ऐसा कलनविधि है जो परिमित समय में बता दे कि पूर्णांक गुणांकों
वाले किसी बहुपद समीकरण, जैसे $3x^2 - 2y +3 = 0$, का पूर्णांकों में हल है या
नहीं। ऐसे प्रश्नों को \emph{डायोफैंटाइन समस्याएँ} कहते हैं। कुछ आरंभिक
सफलताओं के बाद रॉबिन्सन ने मार्टिन डेविस और हिलरी पटनम के साथ काम आरंभ किया,
जो इसी समस्या पर काम कर रहे थे। वे यह दिखाने में सफल हुए कि घातीय
डायोफैंटाइन समस्याएँ (जिनमें अज्ञात राशियाँ घातांक के रूप में भी आ सकती हैं)
अनिर्णेय हैं, और उन्होंने दिखाया कि एक निश्चित अनुमान (जिसे बाद में
``जे.आर.'' कहा गया) से हिल्बर्ट की दसवीं समस्या की अनिर्णेयता निष्पन्न होती
है \citep{DavisPutnamRobinson1961}। रॉबिन्सन पूरे 1960 के दशक में समस्या पर
काम करती रहीं। 1970 में युवा रूसी गणितज्ञ यूरी मतियासेविच ने अंततः जे.आर.
परिकल्पना सिद्ध की। संयुक्त परिणाम को अब
मतियासेविच--रॉबिन्सन--डेविस--पटनम प्रमेय, संक्षेप में एमआरडीपी प्रमेय, कहते
हैं। मतियासेविच और रॉबिन्सन मित्र बने और उन्होंने कई शोधपत्रों पर सहयोग
किया। मतियासेविच को लिखे एक पत्र में रॉबिन्सन ने कहा था, ``वास्तव में मुझे
बहुत प्रसन्नता है कि साथ काम करते हुए (हज़ारों मील की दूरी से) हम स्पष्टतः
उससे अधिक प्रगति कर रहे हैं जितनी हममें से कोई अकेले कर पाता''
\citep[45]{Matijasevich1992}।

रॉबिन्सन अमेरिकन मैथमेटिकल सोसाइटी की पहली महिला अध्यक्ष और नेशनल अकैडमी ऑफ
साइंस के लिए निर्वाचित पहली महिला थीं। ल्यूकीमिया का निदान होने के बाद 30
जुलाई 1985 को 65 वर्ष की आयु में उनका निधन हुआ।

\begin{reading}
रॉबिन्सन के गणितीय शोधपत्र उनकी \textit{Collected Works}
\citep{Robinson1996} में उपलब्ध हैं, जिसमें उनका नेशनल अकैडमी ऑफ साइंसेज़
जीवनीपरक संस्मरण \citep{Feferman1994} भी पुनर्मुद्रित है। रॉबिन्सन की बड़ी
बहन कॉन्स्टन्स रीड ने साक्षात्कारों पर आधारित ``Autobiography of Julia''
\citep{Reid1986} और एक पूर्ण संस्मरण \citep{Reid1996} प्रकाशित किया। जॉर्ज
चिचेरी ने रॉबिन्सन और हिल्बर्ट की दसवीं समस्या पर एक लघु वृत्तचित्र निर्देशित
किया \citep{Csicsery2016}। रॉबिन्सन के साथ यूरी मतियासेविच के सहयोग और उनके
कार्य पर रॉबिन्सन के प्रभाव के संक्षिप्त संस्मरण के लिए
\citep{Matijasevich1992} देखें।
\end{reading}

\end{document}
